%\documentclass{standalone}
%\usepackage{tikz}
%\usetikzlibrary{patterns,calc}
%\begin{document}
\begin{tikzpicture}[scale=1.25]

% Draw the ambient space X as a rectangle
\draw[thick, gray] (-2.9,-2.8) rectangle (5.9,3.2) node[anchor=north east] {$X$};

% Draw the compact set K as a closed irregular shape

\draw[very thick, black] plot [smooth cycle] coordinates {(-0.3,0.1) (-1,0)  (-1,-0.5) (-1.5,-1) (-0.5,-1.3) (0,-0.6) (1,-1.3) (1.4,-0.9) (0.6,0) (0.7,0.8) (-0.1,1)};
\node at (1.05,-1) {$K$};

% Points x_1, x_2, x_3 inside K
\node[name=x1] at (-0.9,-0.85) {};
\node[name=x2] at (0.73,-0.635) {};
\node[name=x3] at (0.2,0.6) {};
\filldraw[black]  (x1) circle (2pt); 
\filldraw[black]  (x2) circle (2pt); 
\filldraw[black]  (x3) circle (2pt); 

% Point z outside K
\node[name=z]  at (3,1.6) {};
\filldraw[black] (z) circle (2pt);

% Open neighborhoods U(x_1), U(x_2), U(x_3) covering K,
% and their counterparts - 
% open neighborhoods V_{x_1}(z),  V_{x_2}(z), V_{x_3}(z)
% such that U(x_i) is disjoint from V_{x_i}(z).
% We draw, shade and label these open neighborhoods
\path[save path=\Uxone] (x1) circle[x radius=1.3, y radius=0.82, rotate=45];
\path[save path=\Vxone] (z)  circle[x radius=1.3, y radius=0.82, rotate=45];

\draw[use path=\Uxone, blue, thick, dashed];
\fill[use path=\Uxone, pattern=north east lines, pattern color=blue, opacity=0.2];
\draw[use path=\Vxone, blue, thick, dashed];
\fill[use path=\Vxone, pattern=north east lines, pattern color=blue, opacity=0.1];
\node[anchor=north, blue] at ($(x1)+(0,-1.2)$) {$U(x_1)$};
\node[anchor=north, blue] at ($(z)+(0,-1.2)$) {$V_{x_1}(z)$};


\path[save path=\Uxtwo] (x2) circle[x radius=1.2, y radius=0.6, rotate=145];
\path[save path=\Vxtwo] (z) circle[x radius=1.2, y radius=0.6, rotate=145];
\draw[use path=\Uxtwo, red, thick, dashed];
\fill[use path=\Uxtwo, pattern=north west lines, pattern color=red, opacity=0.1];
\draw[use path=\Vxtwo, red, thick, dashed];
\fill[use path=\Vxtwo, pattern=north west lines, pattern color=red, opacity=0.1];
\node[anchor=north west, red] at ($(x2)+(1.05,-0.7)$) {$U(x_2)$};
\node[anchor=north west, red] at ($(z)+(1.05,-0.7)$) {$V_{x_2}(z)$};


\path[save path=\Uxthree] (x3) circle[x radius=2.1, y radius=0.55]; 
\path[save path=\Vxthree] (z) circle[x radius=2.1, y radius=0.55]; 
\draw[use path=\Uxthree, green!50!black, thick, dashed] ;
\fill[use path=\Uxthree, pattern=horizontal lines, pattern color=green!50!black, opacity=0.2];
\node[anchor=south east, green!50!black] at ($(x3)+(-1.6,0.3)$) {$U(x_3)$}; 
\draw[use path=\Vxthree, green!50!black, thick, dashed] ;
\fill[use path=\Vxthree, pattern=horizontal lines, pattern color=green!50!black, opacity=0.2];
\node[anchor=south east, green!50!black] at ($(z)+(-1.6,0.3)$) {$V_{x_3}(z)$}; 

% The open neighborhood V(z) is defined as V_{x_1}(z) ∩ V_{x_2}(z) ∩ V_{x_3}(z)
\begin{scope}
	\clip[use path=\Vxone];
	\clip[use path=\Vxtwo];
	\fill[use path=\Vxthree, gray, opacity=0.5];
\end{scope}

% Label the neighborhood V(z)
\node[black,anchor=south] at (z) {$V(z)$};
\end{tikzpicture}
%\end{document}